% chapters/13/sub-cpmk.tex
\begin{subcpmkbox}
Bab ini menuntaskan \SUBCPMK{3.2}. Ketiga indikatornya diperiksa, dan buktinya
berupa level yang dapat dimainkan beserta catatan perbaikan masalah fisika.

\vspace{0.5em}
\textbf{\SUBCPMK{3.2} Membangun level 2D/3D sederhana dengan
\istilah{collider}, \istilah{Rigidbody}, dan kontrol pemain} --- ranah C4,
\CPMK{3}, \CPL{04}.

\vspace{0.5em}
{\setlength{\tabcolsep}{4pt}\footnotesize%
\begin{tabularx}{\textwidth}{@{}C{1.9cm} L{4.0cm} L{2.5cm} Y@{}}
\toprule
\headrow \textbf{Indikator} & \textbf{Rumusan} & \textbf{Dibahas pada} & \textbf{Bukti} \\
\midrule
\ind{a} & Menyusun level menggunakan \istilah{tilemap}/\istilah{greyboxing} &
Bagian~\ref{sec:menyusun-level} &
Level yang tersusun dari petak atau balok abu-abu, beserta tangkapan layarnya
pada laporan praktikum. \\
\addlinespace
\ind{b} & Mengonfigurasi \istilah{Collider} dan \istilah{Rigidbody} sesuai
kebutuhan mekanik &
Bagian~\ref{sec:collider-rigidbody} &
Level yang dapat dimainkan: pemain bergerak, melompat, mendarat di atas
bongkahan, dan terhalang oleh dinding. \\
\addlinespace
\ind{c} & Menganalisis dan memperbaiki masalah fisika &
Bagian~\ref{sec:masalah-fisika} &
Catatan empat masalah fisika (getaran, menembus collider, jatuh tanpa henti,
tersangkut) beserta penyebab dan perbaikannya. \\
\bottomrule
\end{tabularx}}
\end{subcpmkbox}
